\section{One Endpoint, Many Owners}
\label{sec:ch05-one-endpoint-many-owners}

\index{subgraph!is an ordinary service|(}%
A subgraph is your service with four additions. It still serves its own schema
at its own address, it still answers ordinary queries from anyone who asks, and
you can still open it in an IDE and browse it as chapter~\ref{ch:hot-chocolate-zero}
did. Chapter~\ref{ch:do-you-need-this} counted the four as a cost, and they are;
here is what each one buys.

\index{link@\code{"@link}!opts a schema into federation}%
\index{service@\code{\_service}!publishes the schema}%
\index{entities@\code{\_entities}!the entity route}%
The schema applies \code{@link} to itself, which is the opt-in: it says which
version of the federation vocabulary this schema is written in and which
directives it uses. It defines those directives, so that a composer reading the
schema alone knows what they mean. It resolves \code{\_service}, a field whose
one job is to hand back the schema as a string, so the composer can ask the
running service what it now looks like instead of being sent a file. And it
resolves \code{\_entities}, which is the entity route: a root field that takes
keys and returns the objects those keys name.

\index{subgraph!the additions are public}%
None of those four is hidden. They are fields on \code{Query} and directives in
the schema, visible to anybody who introspects the service. That is the fact
chapter~\ref{ch:entities-authorization} is built on.

\subsection{The Schema Nobody Writes}

\index{supergraph!is a composed result}%
Take three subgraph schemas and run them through a composer, and what comes out
is one schema in which \code{Session} carries the fields Sessions declared and
the fields Ratings declared, as though a single service had always had them.
That result is the supergraph. It is not a file anyone maintains, it is not
checked into a repository as a source document, and no team owns it, because it
is derived from the three that are owned.

\index{supergraph!is not a specification term}%
It is also not a word the specification defines. The Apollo Federation subgraph
specification regulates exactly one thing, the subgraph: what it must expose,
what those fields must return, and what the directives mean~\autocite{apollosubgraphspec}.
\enquote{Supergraph} appears in Apollo's own conceptual documentation, under a
heading that introduces it as terminology to learn, defined as \enquote{the
single, federated graph} you get \enquote{when combining multiple GraphQL
APIs}~\autocite{apollofederatedschemas}. That is a useful word and I use it
throughout this book. Which of the two has normative text behind it matters
when two implementations disagree, because the subgraph contract is the only
one of them you can hold anybody to.

\index{router!serves the supergraph}%
The router serves the supergraph, and it is the only address a client has. A
client sends it an ordinary GraphQL operation against an ordinary-looking
schema and gets an ordinary response. Nothing in that response says which
service produced which field, and no part of the protocol offers to tell the
client. That is the design working. It is also
chapter~\ref{ch:blame-routing}'s entire problem, arriving early.

\subsection{Two Different Times, Two Different Failures}

\index{composition!happens before any request}%
The join in figure~\ref{fig:ch05-the-join-moves} needs somebody to have worked
out, in advance, that a session's speaker is answered by asking the Speakers
service for an entity. Working that out is composition, and it happens before a
single request has been served. It reads the three schemas, checks that every field is
reachable from the root through some sequence of services, and produces
something the router can execute against.

\index{composition!fails on somebody else's change}%
Composition fails on schemas, not on data. It fails after everyone has merged,
it fails for the graph rather than for one service, and it can fail because of
a change you did not make. Chapter~\ref{ch:composition} is where you run it,
chapter~\ref{ch:satisfiability} is about the errors that are genuinely hard to
read, and appendix~\ref{app:comperrors} decodes the rest.

\index{execution!happens per request}%
Execution is the other time. A query arrives, the router works out which
services to ask in which order, asks them, and merges what comes back. Its
failures are the ones any distributed system has: a service that is slow, a
service that is down, a field that resolves to null and takes a branch of the
response with it. Part~III is almost entirely
about those.

\index{failure!which time it belongs to}%
Sorting a failure into the right one of those two is, in my experience, most of
debugging a federated graph, and the sort is cheap once you think to make it. A
field that has not worked since the deployment, for anyone, is composition. A
field that works for some requests and not others, or worked an hour ago, is
execution. The two have separate tooling, separate logs and separate people to
ask, so the cost of guessing wrong is an afternoon.
\index{subgraph!is an ordinary service|)}%
